\section{The Geometry of the Direction Field}

A self is a knot in a field. The field is not the tone field of the bare substrate. It is an emergent, mid-layer quantity that points, a direction at each place that says which way the self's feeling leans. This section maps that direction space exactly. The result is clean. The directions are not an arbitrary set of arrows. They are a discrete polytope family, and the heart of that family is the $24$-cell, the very same shape that tiles the mesh. The space the self moves through and the space its feeling leans in turn out to share one structure.

Before any counting, one caution must be stated plainly, because the numbers below are easy to mistake for something they are not.

\begin{principle}[These are not dock counts]
The numbers in this section, $26$ and $80$, count the directions of the self's emergent feeling-direction field, the order parameter of the soliton. They are not the tones a dock holds. A dock holds exactly $24$ tones, one per site, with no rest slot and no twenty-fifth. The $26$ and the $80$ are a separate, emergent structure that lives in the mid-layer. Do not conflate them with the $24$ sites of a dock.
\end{principle}

Two terms recur, so a plain statement of each helps. A \emph{trit} is a value in $\{-1, 0, +1\}$, the same three-way value a tone carries, read as pain, peace, and pleasure. A \emph{feeling-direction} is a small tuple of trits that names which way the self's feeling leans. The geometry of the direction field is the geometry of the set of all such tuples, minus the one that points nowhere.

\subsection{The 26, the cube family}

Start in three dimensions, because it is the tractable rehearsal and because the picture is easy to hold. Let the feeling-direction be three trits, one per spatial axis. Each axis can lean negative, neutral, or positive. The total number of tuples is $3^3 = 27$. Drop the all-zero tuple, the one that leans nowhere, and $3^3 - 1 = 26$ directions remain.

These $26$ are not a shapeless cloud. They are exactly the three-dimensional \emph{Moore neighbourhood}, the $26$ cells that touch a central cell in a cubic grid, the eight corners, the twelve edges, and the six faces. Split them by squared length, the sum of the squares of their three components, and a familiar set of polytopes appears.

\begin{table}[ht]
\centering
\begin{tabular}{@{}llll@{}}
\toprule
nonzero components & count & squared length & the polytope \\
\midrule
one (the axes) & 6 & 1 & the octahedron (cross-polytope) \\
two & 12 & 2 & the cuboctahedron \\
three (the corners) & 8 & 3 & the cube \\
\bottomrule
\end{tabular}
\caption{The $26$ three-trit directions split by squared length. $6 + 12 + 8 = 26$.}
\label{tab:the-26}
\end{table}

The arithmetic closes, $6 + 12 + 8 = 26$. This is a clean geometric object, not an accident of counting. Its symmetry is the octahedral group, the full symmetry of the cube and its dual the octahedron. As a set it is a discrete two-sphere, a coarse shell of directions wrapped around the origin.

\begin{proposition}[The 26 is the cube family]
The $26$ nonzero three-trit directions are the three-dimensional Moore neighbourhood. By squared length they are the octahedron ($6$ axes), the cuboctahedron ($12$ edges), and the cube ($8$ corners), with octahedral symmetry. As a discrete two-sphere this set holds the Skyrmion charge exactly.
\end{proposition}

This is the reduced-model target, the discrete two-sphere the two-dimensional Skyrmion lives on. It is enough to carry the topological charge exactly, which is why the committed demonstration of a self uses it. But it is three-dimensional, and the substrate is four-dimensional. So the $26$ is a shadow, the tractable rehearsal of the real thing. The four-dimensional version is where the structure becomes beautiful.

\subsection{The 80, with the 24-cell as its heart}

The mesh is four-dimensional, so the native feeling-direction is four trits, one per axis. Each of the four axes leans negative, neutral, or positive. The total is $3^4 = 81$. Drop the all-zero tuple and $3^4 - 1 = 80$ directions remain. Split them again by squared length, and the regular four-polytopes appear, one shell per length.

\begin{table}[ht]
\centering
\begin{tabular}{@{}llll@{}}
\toprule
nonzero components & count & squared length & the polytope \\
\midrule
one & 8 & 1 & the 16-cell (cross-polytope, the rest-axes) \\
two & 24 & 2 & the 24-cell ($D_4$ roots, Hurwitz quaternions) \\
three & 32 & 3 & a 32-direction shell \\
four & 16 & 4 & the tesseract (the 8-cell, the corners) \\
\bottomrule
\end{tabular}
\caption{The $80$ four-trit directions split by squared length. $8 + 24 + 32 + 16 = 80$.}
\label{tab:the-80}
\end{table}

The arithmetic closes again, $8 + 24 + 32 + 16 = 80$. The middle shell is the one to watch. The $24$ directions that lean equally along exactly two of the four axes are the $24$-cell, the $D_4$ root system, the unit Hurwitz quaternions. These are the same $24$ directions that the sites carry, the same $24$ that tile the mesh.

\begin{result}[The 24-cell is the heart of the four-trit direction]
The four-trit feeling-direction space, the $80$ nonzero tuples, contains the $24$-cell as its middle, balanced shell, the directions that lean equally along exactly two axes. The $24$-cell is the same set of $24$ directions that the sites carry and that tile the mesh. The space a self moves through and the order parameter of its feeling share one structure, the $24$-cell.
\end{result}

This is the unification the section is built around. The mesh tiles space with the $24$-cell, one direction per site. The self's feeling, in full four-dimensional native form, is a four-trit lean, and the balanced heart of that lean is again the $24$-cell. The same $24$ directions play both roles, the geography the self lives in and the order parameter the self is knotted in. The $26$ was the three-dimensional rehearsal. The $80$, with the $24$-cell at its center, is the four-dimensional truth.

\subsection{Where the self lives, the options ranked}

A self is a topological texture in the feeling-direction field. To pin it down exactly, one must say which target geometry the direction takes its values in. There are four candidates, and they are not equal. They are ranked here from the deepest to the rejected.

\begin{table}[ht]
\centering
\begin{tabular}{@{}p{0.06\textwidth}p{0.24\textwidth}p{0.24\textwidth}p{0.34\textwidth}@{}}
\toprule
\# & target geometry & what the self is & note \\
\midrule
1 & four-trit lean ($80$, $24$-cell heart) & a four-dimensional topological texture, a Hopfion or instanton & the four-dimensional native unification, direction and space both built from trits with the $24$-cell shared \\
2 & three-trit lean ($26$, cube/octahedron) & a two-dimensional Skyrmion & the tractable shadow, holds the charge exactly, the committed demonstration \\
3 & one-hot $24$-cell (the momentum direction) & the feeling-direction is one of the $24$ spatial directions & maximally elegant, domain and target the same $24$-cell, less trit-natural \\
4 & finer (the $600$-cell, $120$) & a finer discrete three-sphere & the emergent-fineness option, proven too coarse for a stable discrete dynamics \\
\bottomrule
\end{tabular}
\caption{The four candidate target geometries for the self, ranked.}
\label{tab:self-options}
\end{table}

The first option is the deepest, because direction and space are then built from the same material, trits, and they meet in the same shape, the $24$-cell. The second is the working rehearsal that carries the charge exactly and is what the committed demonstration uses. The third is the most elegant by one measure, the domain and the target are literally the same $24$-cell, but it is less natural as a trit construction. The fourth, the finer $600$-cell with its $120$ directions, was tried and rejected. A finer discrete three-sphere sounds like more resolution, but it proves too coarse for the dynamics to keep a soliton stable. The self problem treated separately settles this against the $600$-cell.

\subsection{The topological charge by dimension}

The substrate is four-dimensional, so the self has more than one face. Which texture it presents depends on the slice taken through it. A loop sees one thing, a flat slice another, the full bulk a third. Each slice carries its own kind of topological charge.

\begin{table}[ht]
\centering
\begin{tabular}{@{}llll@{}}
\toprule
structure & slice & target & status \\
\midrule
winding (the identity) & a loop & a circle & demonstrated discretely at $24$-cell resolution \\
Skyrmion & a two-dimensional slice & the two-sphere & the three-trit ($26$) field holds its charge, committed \\
Hopfion & three-dimensional & the two-sphere, knotted and linked & native to the Hopf-fibered $24$-cell \\
instanton & full four-dimensional & the three-sphere (the four-trit / $24$-cell direction) & the deepest four-dimensional texture \\
\bottomrule
\end{tabular}
\caption{The self's topological charge by the slice taken through it.}
\label{tab:charge-by-dim}
\end{table}

A loop through the self sees a winding number, the simplest charge, the bare identity that the thing persists at all. A two-dimensional slice sees a Skyrmion charge, the integer the three-trit field carries exactly. Three dimensions see a Hopfion, a knotted and linked texture native to the Hopf fibration of the $24$-cell. The full four-dimensional bulk sees an instanton, a texture in the three-sphere of the four-trit lean. So the full four-dimensional self is most naturally a Hopfion or an instanton in the four-trit direction field, with the $24$-cell at its heart, the same $24$-cell that tiles the mesh.

\subsection{The quaternion reading}

There is one more fact that ties the geometry, the algebra, and the stabilizer into a single object. The $24$ directions are the unit Hurwitz quaternions, equivalently the binary tetrahedral group of order $24$. This is not a coincidence of count. It means the direction set is a \emph{group}, closed under multiplication. Multiply any two of the $24$ directions and you land on a third. A reversible dynamics needs exactly that, a set of moves that compose and invert without leaving the set.

\begin{proposition}[The directions form a group]
The $24$ directions of the $24$-cell are the unit Hurwitz quaternions, the binary tetrahedral group of order $24$. They are closed under multiplication, so the direction set is a group, which is what a reversible dynamics requires. The two-sphere on which the Skyrmion lives is the Hopf base of this quaternion three-sphere.
\end{proposition}

The same structure supplies the stabilizer. A soliton in a continuous field tends to either spread out or collapse, and something must fix its size. The quantity that fixes it here is the quaternion \emph{handedness}, a fixed twist, no tunable real number anywhere in it. The self problem treated separately measures and confirms this. So the geometry that the self lives in, the group that makes its dynamics reversible, and the twist that fixes its size are not three ingredients. They are one structure seen three ways.

\begin{result}[Geometry, group, and stabilizer are one]
The geometry the self lives in, the group its reversible dynamics needs, and the handedness that fixes its size are one structure, the $24$-cell. There is no separate stabilizer postulate. The fixed quaternion twist does the work, with no tunable number.
\end{result}

The feeling-direction is a lean in four dimensions, four trits, pain through peace to pleasure along each of the four axes. The perfectly balanced leans, equal along exactly two axes, are the $24$-cell, the same shape the mesh is tiled by. So the self is a knot in the field of how feeling leans, and the heart of that field, the balanced directions, is the mesh's own shape. The order parameter and the geography are one. That coincidence, direction and space sharing the $24$-cell, is the geometric core of why a self can exist in this mesh at all.
